\section*{Bab \ref{ch:g7:priorities} --- Urutan Operasi Hitung}

\begin{solution}{exo:g7:priorities:1}
$5 + 3 \times 6 = 5 + 18 = 23$.

$(5 + 3) \times 6 = 8 \times 6 = 48$.

$30 - 12 \div 4 = 30 - 3 = 27$.

$(30 - 12) \div 4 = 18 \div 4 = 4.5$.
\end{solution}

\begin{solution}{exo:g7:priorities:2}
$24 - 6 + 2 = 18 + 2 = 20$ (kiri ke kanan).

$24 \div 6 \times 2 = 4 \times 2 = 8$ (kiri ke kanan).

$24 - 6 \times 2 = 24 - 12 = 12$ (perkalian lebih dulu).

$24 \div (6 \times 2) = 24 \div 12 = 2$.
\end{solution}

\begin{solution}{exo:g7:priorities:3}
$7 + 3 \times (10 - 6) = 7 + 3 \times 4 = 7 + 12 = 19$.

$40 - (2 + 3) \times 4 = 40 - 5 \times 4 = 40 - 20 = 20$.

$60 \div \bigl(2 \times (7 - 4)\bigr) = 60 \div (2 \times 3)
= 60 \div 6 = 10$.
\end{solution}

\begin{solution}{exo:g7:priorities:4}
$\dfrac{18 + 6}{8} = \dfrac{24}{8} = 3$; \quad
$\dfrac{30}{7 - 2} = \dfrac{30}{5} = 6$; \quad
$\dfrac{5 \times 8}{4 + 6} = \dfrac{40}{10} = 4$.
\end{solution}

\begin{solution}{exo:g7:priorities:5}
$8 \times 102 = 8 \times (100 + 2) = 800 + 16 = 816$.

$25 \times 99 = 25 \times (100 - 1) = 2500 - 25 = 2475$.

$14 \times 12 + 14 \times 8 = 14 \times (12 + 8) = 14 \times 20 = 280$.

$37 \times 45 - 37 \times 35 = 37 \times (45 - 35) = 37 \times 10 =
370$.
\end{solution}

\begin{solution}{exo:g7:priorities:6}
``\omterm{ex:g1:subtraction:difference}{Selisih} dari $50$ dan \omterm{def:g2:mult:def}{hasil kali} $6$ dengan $7$'':
$50 - 6 \times 7 = 50 - 42 = 8$ (tidak perlu tanda kurung, urutan
pengerjaannya sudah bekerja).

``\omterm{def:g2:mult:def}{Hasil kali} dari \omterm{def:g1:addition:def}{jumlah} $4$ dan $5$ dengan \omterm{ex:g1:subtraction:difference}{selisih} $10$ dan
$8$'': $(4 + 5) \times (10 - 8) = 9 \times 2 = 18$.
\end{solution}

\begin{solution}{exo:g7:priorities:7}
$3 \times 2.50 + 2 \times 1.20 = 7.50 + 2.40 = 9.90$ euro.
\end{solution}

\begin{solution}{exo:g7:priorities:8}
(a) tanpa kurung: $4 + 2 \times 5 - 1 = 4 + 10 - 1 = 13$.

(b) $(4 + 2) \times 5 - 1 = 30 - 1 = 29$.

(c) $(4 + 2) \times (5 - 1) = 6 \times 4 = 24$.
\end{solution}

\begin{solution}{exo:g7:priorities:9}
$k = 1$: $12 + 4 \times 1 = 16$ dan $16 \times 1 = 16$ --- keduanya
cocok (itulah sebabnya $k = 1$ adalah uji yang buruk!). $k = 2$:
$12 + 4 \times 2 = 20$ tetapi $16 \times 2 = 32$: berbeda. Jadi kedua
bentuk itu \emph{tidak} sama secara umum. Kesalahan
$12 + 4 \times k = 16 \times k$ datang dari menjumlahkan $12 + 4$
lebih dulu, dengan mengabaikan urutan pengerjaan $\times$.
\end{solution}

\begin{solution}{exo:g7:priorities:10}
\emph{1.} Harga rombongan: $12 + 6 \times n$. Untuk $n = 5$:
$12 + 30 = 42$ euro.

\emph{2.} Harga satuan: $9n$. Bandingkan: $n = 3$: $30$ lawan $27$
(satuan lebih murah); $n = 4$: $36$ lawan $36$ (sama); $n = 5$: $42$
lawan $45$ (rombongan lebih murah). Mulai dari $5$ orang, tarif
rombongan menang (setiap orang tambahan berharga $6$ dan bukan $9$,
jadi jaraknya hanya bertambah).
\end{solution}

\begin{solution}{exo:g7:priorities:11}
Jawaban yang mungkin (ada banyak):
\[
24 = 1 \times 2 \times 3 \times 4, \qquad
10 = 1 + 2 + 3 + 4, \qquad
1 = (4 - 3) \times (2 - 1).
\]
\end{solution}

\begin{solution}{pb:g7:priorities:1}
\textbf{1.} $7 \times 102 = 7 \times (100 + 2) = 700 + 14 =
714$; \quad $9 \times 98 = 9 \times (100 - 2) = 900 - 18 = 882$.

\textbf{2.} $17 \times 13 + 17 \times 7 = 17 \times (13 + 7) =
17 \times 20 = 340$; \quad $45 \times 101 - 45 = 45 \times (101
- 1) = 45 \times 100 = 4\,500$.

\textbf{3.} $5 \times 87 = 870 \div 2 = 435$. Kesahihannya:
$5 = 10 \div 2$, jadi mengalikan dengan $5$ berarti mengalikan dengan
$10$ lalu membagi dengan $2$ (urutan $\times$ dan $\div$ yang
setingkat, dari kiri ke kanan, mengerjakan sisanya).

\textbf{4.} $25 \times 32 = 3\,200 \div 4 = 800$, karena
$25 = 100 \div 4$. Dan
$99 \times 47 = (100 - 1) \times 47 = 4\,700 - 47 = 4\,653$:
sifat distributif terhadap \omterm{ex:g1:subtraction:difference}{selisih}
(\cref{thm:g7:priorities:distributivity}).

\textbf{5.} $8 \times (5 + 3) = 64$ dan $8 \times 5 + 8 \times 3
= 40 + 24 = 64$: keduanya cocok --- itulah sifat distributif. Tetapi
$8 \times (5 \times 3) = 8 \times 15 = 120$, sedangkan
$(8 \times 5) \times (8 \times 3) = 40 \times 24 = 960$: tidak
cocok. Faktornya menyebar ke atas \emph{suku sebuah \omterm{def:g1:addition:def}{jumlah}}, bukan ke
atas faktor sebuah \omterm{def:g2:mult:def}{hasil kali}.

\textbf{6.}
\begin{center}
\renewcommand{\arraystretch}{1.15}
\begin{tabular}{ccl}
$13$ & $24$ & disimpan ($13$ \omterm{def:g3:numbers:evenodd}{ganjil}) \\
$6$ & $48$ & dicoret ($6$ \omterm{def:g3:numbers:evenodd}{genap}) \\
$3$ & $96$ & disimpan \\
$1$ & $192$ & disimpan \\
\end{tabular}
\end{center}
\omterm{def:g1:addition:def}{Jumlah} bilangan kanan yang selamat: $24 + 96 + 192 = 312$, dan
memang $13 \times 24 = 312$.

\textbf{7.} Membagi dua $21$: $21, 10, 5, 2, 1$; melipatduakan $17$:
$17, 34, 68, 136, 272$. Yang kirinya \omterm{def:g3:numbers:evenodd}{genap}, $10$ dan $2$, dicoret;
yang selamat $17 + 68 + 272 = 357$. Periksa:
$21 \times 17 = 357$.

\textbf{8.} Membagi dua $16$: $16, 8, 4, 2, 1$; melipatduakan $25$:
$25, 50, 100, 200, 400$. Setiap bilangan kirinya kecuali $1$ yang
terakhir adalah \omterm{def:g3:numbers:evenodd}{genap}: hanya satu baris yang selamat, dan jawabannya
$400$. Kolom kiri yang tersusun dari lipatan dua murni hanya menyimpan
baris terakhirnya --- caranya menyusut menjadi melipatduakan saja.

\textbf{9.} $13 = 12 + 1$, jadi menurut sifat distributif
$13 \times 24 = 12 \times 24 + 1 \times 24$; dan
$12 \times 24 = 6 \times 48$ karena membagi dua satu faktor sambil
melipatduakan faktor yang lain membiarkan \omterm{def:g2:mult:def}{hasil kalinya} tetap
($12 \times 24 = 6 \times 2 \times 24 = 6 \times 48$). Bilangan kiri
$13$ yang \omterm{def:g3:numbers:evenodd}{ganjil} memisahkan satu salinan bilangan kanannya: baris
pertama menabung $24$.

\textbf{10.} $6 \times 48 = 3 \times 96$ (bagi dua--lipat dua, tanpa
menabung: $6$ \omterm{def:g3:numbers:evenodd}{genap}). $3 \times 96 = 2 \times 96 + 96 =
1 \times 192 + 96$: si \omterm{def:g3:numbers:evenodd}{ganjil} $3$ menabung $96$. Akhirnya
$1 \times 192 = 192$: si \omterm{def:g3:numbers:evenodd}{ganjil} $1$ menabung $192$ dan tabelnya
berakhir. Seluruh tabungan: $24 + 96 + 192 = 312$. Sebuah baris
menabung bilangan kanannya tepat ketika bilangan kirinya \omterm{def:g3:numbers:evenodd}{ganjil} ---
jadi resep ``simpan baris yang kirinya \omterm{def:g3:numbers:evenodd}{ganjil} lalu jumlahkan'' adalah
persis pembukuan sifat distributif.

\textbf{11.} $1 + 4 + 8 = 13$. Untuk $21$, baris yang selamat adalah
baris $17$, $68 = 4 \times 17$ dan $272 = 16 \times 17$: yaitu
lipatan dua $1 + 4 + 16 = 21$.

\textbf{12.} $25 = 16 + 8 + 1$ dan $42 = 32 + 8 + 2$.

\textbf{13.} Ambillah bilangan lipatan dua terbesar yang tidak
melebihi sasarannya lalu kurangkan; yang tersisa lebih kecil, dan
lebih kecil daripada bilangan lipatan dua yang baru dipakai (kalau
tidak, yang lebih besar akan muat). Kalau diulang, sisanya jelas
mengecil dan pasti mencapai $0$: prosesnya berhenti, dan lipatan dua
yang dipakai semuanya berbeda. Untuk $100$: $64$ muat (tersisa $36$),
$32$ muat (tersisa $4$), $4$ muat (tersisa $0$): $100 = 64 + 32 + 4$.

\textbf{14.} Tabel lipatan dua dari $35$: $1 \to 35$, $2 \to 70$,
$4 \to 140$, $8 \to 280$, $16 \to 560$. Karena
$19 = 16 + 2 + 1$, pilihlah baris itu:
$560 + 70 + 35 = 665$. Periksa: $19 \times 35 = 665$.

\textbf{15.} Ketiganya sama seperti milik juru tulis dan petani tadi:
$24 + 96 + 192$ --- lipatan dua dari $24$ yang dipilih oleh
$13 = 1 + 4 + 8$. Menulis bilangan sebagai \omterm{def:g1:addition:def}{jumlah} lipatan dua tepat
sama dengan menulisnya dengan angka $0$ dan $1$; ceritanya yang
lengkap dituturkan pada \cref{pb:g8:powers:1}.

\end{solution}
